% Part: normal-modal-logic
% Chapter: axioms-systems
% Section: derived-rules
%READERNOTE{SOL6-A157: replacement-rule च्या संक्षिप्त नोंदीत “नवीन सूत्र for जुने सूत्र” असा क्रम वापरला आहे. म्हणून C(A) वरून C(B) मिळवताना label B for A आहे; हा क्रम खालील double-negation examples मधील p for not-not-p शी जुळतो.}

\documentclass[../../../include/open-logic-section]{subfiles}

\begin{document}

\olfileid{nml}{prf}{der}

\olsection{व्युत्पन्न नियम}

!!{derivation}s शोधणे आणि लिहिणे कठीण, किचकट
आणि पुनरावृत्तीचे काम आहे. उदाहरणार्थ, $!A \lif !B$
पासून $\Box !A \lif \Box !B$ कडे जायचे असेल,
म्हणजे नियम~\RK लावायचा असेल, तर आधी \Nec लावावा लागतो,
मग~\Ax{K} चा योग्य दृष्टांत लिहावा लागतो आणि नंतर~\MP
लावावा लागतो. तसेच $!A \lif !B$ आणि $!B \lif !C$
पासून $!A \lif !C$ कडे जाण्यासाठी पुढील (लांब)
सर्वतः-सत्य दृष्टांत लिहावा लागतो:
\[
(!A \lif !B) \lif ((!B \lif !C) \lif (!A \lif !C))
\]
आणि \MP{} दोनदा लावावा लागतो. अनेकदा एखादे उप-!!{formula}
त्याच्याशी समतुल्य असलेल्या सूत्राने बदलायचे असते;
उदाहरणार्थ, \iftag{prvDiamond}{$\Diamond
  !A$ च्या जागी $\lnot\Box\lnot !A$}{$\Box !A$ च्या जागी
  $\lnot\Diamond\lnot !A$}, किंवा
$\lnot\lnot !A$ च्या जागी $!A$. म्हणून संपूर्ण
!!{derivation} लिहिण्याऐवजी मध्यवर्ती पायऱ्या
!!{derivable} का आहेत, एवढेच नोंदवणे सोयीचे आहे.
त्यासाठी !!{derivability} विषयी काही तथ्ये गोळा करू.

\begin{prop}
  $\Log{K} \Proves !A_1$, \dots, $\Log{K} \Proves !A_n$
  असेल आणि विधानीय तर्कशास्त्रानुसार $!A_1$, \dots,
  $!A_n$ पासून $!B$ निष्पन्न होत असेल, तर
  $\Log{K} \Proves !B$.
\end{prop}

\begin{proof}
  विधानीय तर्कशास्त्रानुसार $!A_1$, \dots,~$!A_n$
  पासून $!B$ निष्पन्न होत असेल, तर
  \[
  !A_1 \lif (!A_2 \lif \cdots (!A_n \lif !B)\dots)
  \]
  हा सर्वतः-सत्य दृष्टांत असतो. \MP{} $n$ वेळा
  लावल्याने $!B$ ची !!{derivation} मिळते.
\end{proof}

या प्रतिज्ञेचा वापर~\PL ने दर्शवू.

\begin{prop}
  $\Log{K} \Proves !A_1 \lif (!A_2 \lif \cdots (!A_{n-1} \lif
  !A_n)\dots)$ असेल, तर $\Log{K} \Proves \Box !A_1 \lif (\Box !A_2 \lif
  \cdots (\Box !A_{n-1} \lif \Box !A_n)\dots)$.
\end{prop}

\begin{proof}
  \olref[nor]{prop:rk} च्या सिद्धतेप्रमाणेच
  $n$ वर विगमन करा.
\end{proof}

या प्रतिज्ञेचा वापर~\RK ने दर्शवू. या परिणामांमुळे
!!{derivability} विषयीची विधाने किती सोप्या
रीतीने सिद्ध करता येतात, ते पाहू.

\begin{prop}
  $\Log{K} \Proves (\Box!A \land \Box!B) \lif \Box (!A \land !B)$
\end{prop}

\begin{proof}
  \begin{derivation}
    1. & $\Log{K} \Proves !A \lif (!B \lif (!A \land !B))$ & \Taut \\
    2. & $\Log{K} \Proves \Box!A \lif (\Box !B \lif \Box(!A \land !B))$ & \RK, 1\\
    3. & $\Log{K} \Proves (\Box!A \land \Box!B) \lif \Box(!A \land !B)$ & \PL, 2
  \end{derivation}
\end{proof}

\begin{prop}\ollabel{prop:rewriting}
  $\Log{K} \Proves !A \liff !B$ आणि $\Log{K} \Proves
  \Subst{!C}{!A}{q}$ असेल, तर $\Log{K} \Proves \Subst{!C}{!B}{q}$.
\end{prop}

\begin{proof}
  सराव.
\end{proof}

\begin{prob}
  $!C$ च्या जटिलतेवर विगमन करून
  \olref[nml][prf][der]{prop:rewriting} सिद्ध करा:
  $\Log{K} \Proves !A \liff !B$ असेल, तर
  $\Log{K} \Proves \Subst{!C}{!A}{q} \liff \Subst{!C}{!B}{q}$
  हे दाखवा.
\end{prob}

$\Diamond$ चे $\Box$ मध्ये (किंवा उलट) रूपांतर करायचे
असेल, किंवा !!a{formula} मधील दुहेरी निषेध काढायचा
असेल, तेव्हा ही प्रतिज्ञा विशेष उपयोगी ठरते.
पुढे $!C(!A)$ च्या जागी !!a{formula}~$!C(!B)$ लिहिताना
\olref{prop:rewriting} चा वापर ``$!B$ for $!A$''
असा दाखवू. दुसऱ्या शब्दांत, ``$!B$ for $!A$''
याचा संक्षिप्त अर्थ असा:
  \begin{derivation}
    & $\Proves !C(!A)$ \\
    & $\Proves !A \liff !B$\\
    & $\Proves !C(!B)$ & \olref{prop:rewriting} वरून
  \end{derivation}
उदाहरणार्थ:

\begin{prop}
  $\Log{K} \Proves \lnot\Box p \lif \Diamond \lnot p$
\end{prop}

\begin{proof}
  \iftag{prvDiamond}{% Diamond is primitive
  \begin{derivation}
    1. & $\Log{K} \Proves \Diamond \lnot p \liff \lnot\Box\lnot\lnot p$ &
    \Dual\\
    2. & $\Log{K} \Proves \lnot \Box \lnot\lnot p \lif \Diamond \lnot p$ & \PL, 1\\
    3. & $\Log{K} \Proves \lnot\Box p \lif \Diamond \lnot p$ & $p$ for $\lnot\lnot p$\\
  \end{derivation}
  }{% Diamond is defined
  \begin{derivation}
    1. & $\Log{K} \Proves \lnot \Box \lnot\lnot p \lif \Diamond \lnot p$ & \Taut\\
    2. & $\Log{K} \Proves \lnot\Box p \lif \Diamond \lnot p$ & $p$ for $\lnot\lnot p$\\
  \end{derivation}
  ओळ~$1$ वरील सूत्र $\lnot q \lif \lnot q$ चा
  दृष्टांत आहे; येथे $q$ च्या जागी $\Box\lnot\lnot p$
  ठेवले आहे. $\Diamond$ ची व्याख्या~$\lnot\Box\lnot$
  असल्याने $\Diamond\lnot p$ आणि
  $\lnot\Box\lnot\lnot p$ ही एकच सूत्रे आहेत.}
\end{proof}

वरील !!{derivation} मधील अंतिम पायरी
``$p$ for $\lnot\lnot p$'' ही पुढील पायऱ्यांचे
संक्षिप्त रूप आहे:
  \begin{derivation}
    & $\Log{K} \Proves \lnot \Box \lnot\lnot p \lif \Diamond \lnot
    p$\\
    & $\Log{K} \Proves \lnot\lnot p \liff p$ & \Taut\\
    & $\Log{K} \Proves \lnot\Box p \lif \Diamond \lnot p$ & \olref{prop:rewriting} वरून
  \end{derivation}
\olref{prop:rewriting} मधील $!C(q)$, $!A$ आणि $!B$
यांच्या जागी येथे अनुक्रमे
$\lnot \Box q \lif \Diamond \lnot p$,
$\lnot\lnot p$ आणि~$p$ आहेत.

!!a{formula} मध्ये $\lnot\Diamond !A$ हे उप-!!{formula}
असेल, तर \olref{prop:rewriting} वापरून त्याच्या जागी
$\Box\lnot !A$ लिहिता येते, कारण
$\Log{K} \Proves \lnot\Diamond !A \liff \Box\lnot !A$.
हा आणि यांसारखे बदल थोडक्यात
``$\Box\lnot$ for $\lnot\Diamond$'' असे दर्शवू.

पुढील प्रतिज्ञेमुळे !!{derivability} विषयीचे परिणाम
रूपबंधात्मक पद्धतीने सिद्ध करता येतात. उदाहरणार्थ,
आधीची प्रतिज्ञा केवळ $\Log{K} \Proves \lnot\Box p \lif
\Diamond \lnot p$ एवढेच दाखवत नाही; कोणत्याही~$!A$
साठी $\Log{K} \Proves \lnot\Box !A \lif \Diamond
\lnot !A$ हेही दाखवते.

\begin{prop}
  $!A$ हा $!B$ चा आदेशन-दृष्टांत असेल आणि
  $\Log{K} \Proves !B$ असेल, तर $\Log{K} \Proves !A$.
\end{prop}

\begin{proof}
  संपूर्ण !!{derivation} वर आदेशन लावल्यानेही योग्य
  !!{derivation} मिळते, हे तपासणे किचकट पण नेहमीचे
  काम आहे ($!B$ च्या !!{derivation} च्या लांबीवर
  विगमन करा). विशेषतः, सर्वतः-सत्य दृष्टांतांचे
  आदेशन-दृष्टांतही सर्वतः-सत्य दृष्टांत असतात;
  \iftag{prvDiamond}{\Dual{} आणि~}{}\Ax{K}
  यांच्या दृष्टांतांचे आदेशन-दृष्टांतही
  \iftag{prvDiamond}{\Dual{} आणि~}{}\Ax{K}
  यांचेच दृष्टांत असतात; आणि आधारसूत्रांमध्ये
  व निष्कर्षात !!{propositional variable}s च्या
  जागी !!{formula}s ठेवले तरी \MP{} व~\Nec{}
  यांचा वापर योग्यच राहतो.
\end{proof}


\end{document}
